%! Author = denis
%! Date = 11/08/2026

\begin{frame}{Identidades lógicas: Lei distributiva}
    
    \begin{block}{Definição}
        Distribuir uma variável sobre uma expressão lógica,
        possibilitando expandir ou fatorar expressões booleanas.
    \end{block}

    \begin{columns}[T]

        % Coluna AND sobre OR
        \begin{column}{0.5\textwidth}
            \begin{block}{AND sobre OR}
                \[
                    \boxed{A(B+C)=AB+AC}
                \]
                \[
                    A(B+C)
                \]
                \[
                    = A\cdot B + A\cdot C
                \]
            \end{block}
        \end{column}

        % Coluna OR sobre AND
        \begin{column}{0.5\textwidth}
            \begin{block}{OR sobre AND}
                \[
                    \boxed{A+BC=(A+B)(A+C)}
                \]
                \[
                    A+BC
                \]
                \[
                    = (A+B)(A+C)
                \]
            \end{block}
        \end{column}

    \end{columns}

\end{frame}